\chapter{Model and training}
\label{ch:methods}

\section{Architecture}

\begin{figure}[t]
\centering
\begin{tikzpicture}[font=\scriptsize, node distance=2.5mm and 3.2mm,
    box/.style={draw, rounded corners=1.5pt, minimum height=6mm, inner sep=2.5pt, align=center},
    acc/.style={box, fill=blue!10, draw=blue!60!black, thick},
    arr/.style={-{Latex[length=1.6mm]}, thin}]
  \node[box] (in) {input spikes $S_t$\\(embedding, L1, LIF)};
  \node[box, right=of in] (qkv) {Q, K, V:\\Linear $\to$ L1 $\to$ LIF};
  \node[acc, right=of qkv] (att) {\csq{} attention\\$\sum_j (q_i\!\cdot\!k_j)^2\, v_j$};
  \node[box, right=of att] (hn) {head L1 $\to$ LIF\\$\to$ proj (L1, LIF)};
  \node[box, below=7mm of hn] (add1) {$x \mathrel{+}= $ spikes};
  \node[box, left=of add1] (mlp) {MLP $256\!\to\!1024\!\to\!256$\\(Linear $\to$ L1 $\to$ LIF)$\times$2};
  \node[box, left=of mlp] (add2) {$x \mathrel{+}= $ spikes};
  \node[box, left=of add2] (out) {mean over $T$\\$\to$ linear head};
  \draw[arr] (in) -- (qkv); \draw[arr] (qkv) -- (att); \draw[arr] (att) -- (hn);
  \draw[arr] (hn) -- (add1); \draw[arr] (add1) -- (mlp); \draw[arr] (mlp) -- (add2); \draw[arr] (add2) -- (out);
  \node[draw, dashed, rounded corners, fit=(qkv)(att)(hn)(add1)(mlp)(add2), inner sep=3mm,
        label={[font=\scriptsize]above:spiking block, $\times 6$, all $T$ time steps}] {};
\end{tikzpicture}
\caption[Architecture of the spiking transformer]{Architecture of the spiking transformer. Inside a block every
arrow carries binary spikes except the residual stream $x$, which holds small integers (counts of spikes added so
far).}
\label{fig:arch}
\end{figure}

The spiking transformer (Figure~\ref{fig:arch}) has the same shape as a standard pre-normalised transformer: 6
blocks of width 256 with 8 attention heads of 32 channels, an MLP of width 1{,}024, a context of 256 tokens and 4.89
million parameters. The differences are in what flows between the layers:

\begin{itemize}
  \item \textbf{Input.} The token and position embeddings are normalised and fed as a constant input current to a
        layer of LIF neurons for $T$ time steps, producing $T$ binary spike vectors per token.
  \item \textbf{Every linear layer is followed by a normalisation and a LIF neuron}, so every layer outputs spikes.
        The normalisation is an \emph{L1} layer norm: it subtracts the mean and divides by the mean absolute
        deviation (MAD) instead of the standard deviation, which avoids a square root.
  \item \textbf{Attention} uses spike counts (Section~\ref{sec:csq}) followed by a per-head L1 normalisation and a
        LIF neuron, then a projection layer.
  \item \textbf{The residual stream} adds the spikes of the attention and MLP outputs to $x$. After block $k$ it
        holds integers from 0 to $2k+3$, at most 13.
  \item \textbf{Read-out.} The residual stream is averaged over the $T$ time steps (a firing rate) and a linear
        head produces the next-token distribution.
\end{itemize}

The non-spiking \emph{teacher} used for distillation is a transformer of exactly the same shape with standard layer
norms, softmax attention and GELU activations (4.87 million parameters).

\section{Spike-count attention and \csq{} attention}
\label{sec:csq}

In spiking attention the score of key $j$ for query $i$ is the overlap count
$c_{ij} = q_i \cdot k_j = \mathrm{popcount}(q_i \wedge k_j)$, an integer from 0 to 32 per head and time step.
Measurements showed that this attention is much broader than the teacher's: rows of normalised counts had a
normalised entropy (the ``spread'') of about 0.8, where 1 means uniform.

\paragraph{Count$^2$ attention.} The rule proposed here squares the counts:
\begin{equation}
  a_i = \sum_{j \le i} c_{ij}^{\,2}\; v_j .
\end{equation}
A key that shares four spikes with the query now counts sixteen times as much as one that shares a single spike,
instead of four times, so strong matches dominate, as with a softmax, but weak matches still contribute a little.
The rule has no parameters and no discontinuities. In hardware the square of a count from 0 to 32 is a lookup in a
33-entry table, and pairs with a count of zero, about three quarters of all pairs in the final model, are skipped.

\paragraph{Count thresholds (explored, not used).} Before \csq{} attention, two threshold-based alternatives were
tried: keeping a count only if it reaches a threshold $\theta$ (a soft winner-take-all), either fixed or learned per
head with a sigmoid surrogate gradient. They are described with their results in Section~\ref{sec:attn}, because
the way they fail is instructive: counts are integers, so every $\theta$ between 1 and 2 has exactly the same effect,
and a $\theta$ that reaches an integer changes a whole head's attention at once.

\section{Training recipe}
\label{sec:recipe}

\paragraph{Distillation.} The training loss combines four terms:
\begin{equation}
  \mathcal{L} = \lambda_{\mathrm{CE}}\, \mathrm{CE} + (1-\lambda_{\mathrm{CE}})\, \mathrm{KL}(p_{\mathrm{teacher}}
  \,\|\, p_{\mathrm{student}}) + \lambda_{\mathrm{att}}\, \mathrm{KL}(\hat{a}_{\mathrm{student}} \,\|\,
  a_{\mathrm{teacher}}) + \lambda_{\mathrm{spread}}\, H(\hat{a}_{\mathrm{student}}) .
\end{equation}
The first two terms distil the output distribution. The third compares the student's attention rows, normalised to
sum to one, with the teacher's softmax attention on 64 randomly sampled query rows per batch in blocks 1--5, head by
head. It uses the \emph{reverse} KL direction: the forward direction made spiking attention flatter, because it
punishes any key the teacher looks at slightly but the student ignores.

\paragraph{Spread penalty.} The fourth term is the normalised entropy of the student's attention rows,
$H/\log(\text{number of visible keys})$. It pushes attention to concentrate, which the distillation loss alone did
not achieve.

\paragraph{Homeostasis.} Some neurons stop firing during training and never recover, because their surrogate
gradient becomes tiny. A homeostatic rule \citep{turrigiano2004homeostatic} nudges the normalisation bias of every
neuron that fires in fewer than 1\% of time steps up, and of every neuron that fires in more than 50\% down, with a
small rate and a cap on the total shift. Without the cap, the rule fought the learned biases and caused gradient
spikes.

\paragraph{Spike dropout.} During training, each block's attention and MLP output spikes are dropped with
probability $p$ before they are added to the residual stream, and kept spikes are scaled by $1/(1-p)$ so the
expected residual is unchanged \citep{srivastava2014dropout}. At inference nothing is dropped and all spikes are
binary.

\paragraph{Schedules.} Count$^2$ attention and the spread penalty are blended in over the first 2{,}000 steps. From
step 2{,}000 to 30{,}000 the output-KL weight, the attention-KL weight and the spread penalty decrease linearly to
zero while the sharpness $\alpha$ of the surrogate gradient increases from 2 to 4; after that training continues on
cross-entropy alone. Every change is linear: switching $\alpha$ abruptly from 2 to 1 in one step, in an early
experiment, raised the validation loss from 1.87 to 2.70, because the optimiser's running estimates of the gradient
scale no longer matched.

\paragraph{Optimisation.} AdamW \citep{loshchilov2019adamw} with learning rate $7 \times 10^{-4}$, 500 warm-up steps
and cosine decay, batches of 64 sequences of 256 tokens, bf16 arithmetic, and a per-block gradient clipping that
limits each block's gradient norm to three times its running average. All hyperparameters are listed in
Appendix~\ref{app:hp}.

\section{Evaluation}
\label{sec:eval}

\paragraph{Prediction.} Bits per token on the full test split, scored with a sliding window of 256 tokens moved by
128, so that every token is predicted with at least 128 tokens of context.

\paragraph{Generated tunes.} 500 tunes are sampled at temperature 1, each starting from the beginning-of-tune token.
Four automatic checks are applied:
\begin{itemize}
  \item \textbf{Correct bars}: the share of bars whose note durations add up to the meter (first and last bars are
        excluded, since pick-up bars are allowed to be short).
  \item \textbf{Header}: exactly one meter followed by exactly one mode at the start.
  \item \textbf{Balanced repeats}: as many closing as opening repeat signs (one missing opening sign at the start
        is allowed) and as many second as first endings.
  \item \textbf{Ends on C}: the last note is the tonic.
\end{itemize}
Real test tunes score 89.9\%, 100\%, 80.3\% and 62.8\% on these checks; the repeat check is stricter than real
practice. With 500 tunes the checks vary by about $\pm$2--3 percentage points between samples.

\paragraph{Baselines.} The ANN teacher (10{,}000 steps) and a folk-rnn-style LSTM with three layers of 512 units and
one-hot input (5.60 million parameters, 30{,}000 steps), both trained on the same split.
